
\documentclass{article}
\usepackage{colortbl}
\usepackage{makecell}
\usepackage{multirow}
\usepackage{supertabular}

\begin{document}

\newcounter{utterance}

\twocolumn

{ \footnotesize  \setcounter{utterance}{1}
\setlength{\tabcolsep}{0pt}
\begin{supertabular}{c@{$\;$}|p{.15\linewidth}@{}p{.15\linewidth}p{.15\linewidth}p{.15\linewidth}p{.15\linewidth}p{.15\linewidth}}

    \# & $\;$A & \multicolumn{4}{c}{Game Master} & $\;\:$B\\
    \hline 

    \theutterance \stepcounter{utterance}  

    & & \multicolumn{4}{p{0.6\linewidth}}{\cellcolor[rgb]{0.9,0.9,0.9}{%
	\makecell[{{p{\linewidth}}}]{% 
	  \tt {\tiny [A$\langle$GM]}  
	 Sono il tuo game master, e stai giocando un gioco collaborativo con la seguente griglia come tavolo da gioco:\\ \tt \\ \tt ```\\ \tt  1234567\\ \tt ╔═══════╗\\ \tt ║◌◌◌◌L◌◌║ 1\\ \tt ║P◌◌◌◌◌◌║ 2\\ \tt ║◌C◌◌◌┌─╢ 3\\ \tt ╟──┐◌◌│◌║ 4\\ \tt ║◌◌├──┼─╢ 5\\ \tt ║◌◌│◌◌│◌║ 6\\ \tt ║◌◌│◌◌│◌║ 7\\ \tt ╚══╧══╧═╝\\ \tt \\ \tt ```\\ \tt \\ \tt * Il bordo superiore mostra le coordinate x che aumentano verso destra, e il bordo destro le coordinate y che aumentano verso il basso.\\ \tt * I seguenti oggetti sono posizionati casualmente sulla tua griglia: 'C', 'L', 'P'.\\ \tt \\ \tt \\ \tt L'altro giocatore vede una variante del tavolo da gioco, dove gli oggetti sono posizionati in diverse posizioni casuali. Non puoi vedere il tavolo dell'altro giocatore, e loro non possono vedere il tuo.\\ \tt \\ \tt \\ \tt **Obiettivo:**\\ \tt \\ \tt Entrambi i giocatori devono spostare gli oggetti sui rispettivi sfondi in modo che oggetti identici finiscano alle stesse coordinate. Devi comunicare con l'altro giocatore per concordare una configurazione obiettivo comune.\\ \tt \\ \tt \\ \tt **Regole:**\\ \tt \\ \tt * In ogni turno, puoi inviare esattamente uno dei seguenti due comandi:\\ \tt  1. `COMUNICA: <MESSAGE>`: per inviare un messaggio (tutto fino alla prossima interruzione di riga) all'altro giocatore. Lo inoltrerò al tuo partner.\\ \tt 2. `SPOSTA: <OBJECT>, (<X>, <Y>)`: per spostare un oggetto in una nuova posizione, dove <X> è la colonna e <Y> è la riga. Ti farò sapere se la tua mossa è andata a buon fine o no.\\ \tt * Se non rispetti il formato, o invii diversi comandi contemporaneamente, devo penalizzarti.\\ \tt * Se entrambi i giocatori accumulano più di 8 penalità, perdete entrambi il gioco.\\ \tt * È vitale che comunichi con l'altro giocatore riguardo al tuo stato obiettivo! L'*unico* modo in cui puoi trasmettere la tua strategia all'altro giocatore è usando il comando `COMUNICA: <MESSAGE>`!\\ \tt \\ \tt \\ \tt **Spostamento Oggetti**\\ \tt \\ \tt * Puoi spostare oggetti solo in celle entro i limiti della griglia. La cella di destinazione deve essere vuota, cioè deve contenere solo il simbolo '◌'.\\ \tt * Se provi a spostare un oggetto in un posto che non è vuoto, o provi a spostarlo fuori dalla griglia, devo penalizzarti. Hai un altro tentativo.\\ \tt * Prima di fare una mossa, controlla due volte che il posto di destinazione sia vuoto, e non contenga alcuna lettera, cornice, o linea!\\ \tt \\ \tt \\ \tt **Fine del Gioco**\\ \tt \\ \tt Se pensi di aver raggiunto l'obiettivo di allineare tutti gli oggetti, puoi chiedere all'altro giocatore di finire il gioco inviando COMUNICA: finito?. Se l'altro giocatore ti chiede di finire il gioco, e rispondi con COMUNICA: finito!, the game will end.\\ \tt \\ \tt Entrambi i giocatori vincono se il gioco finisce entro 12 round, dove un round è definito come due giocatori che inviano ciascuno un comando valido.\\ \tt \\ \tt **Punteggio:**\\ \tt \\ \tt Più vicini sono gli oggetti identici in entrambi i tavoli da gioco, più punti ottieni. Le penalità riducono i tuoi punti. Riesci a battere il record?\\ \tt \\ \tt \\ \tt Per favore invia un messaggio all'altro giocatore per iniziare il gioco! 
	  } 
	   } 
	   } 
	 & \\ 
 

    \theutterance \stepcounter{utterance}  

    & \multicolumn{4}{p{0.6\linewidth}}{\cellcolor[rgb]{0.8,1,0.9}{%
	\makecell[{{p{\linewidth}}}]{% 
	  \tt {\tiny [A$\rangle$GM]}  
	 COMUNICA: Ciao! Il mio tavolo ha: P in (1,2), C in (2,3), L in (5,1). Sto cercando di allineare tutti gli oggetti alla stessa posizione del tuo tavolo. Per favore, fammi sapere dove hai posizionato P, C e L, o suggeriscimi una configurazione obiettivo comune (es: "Mettiamo tutti al centro: P,C,L in (4,4)"). Possiamo farlo in 3 mosse! 
	  } 
	   } 
	   } 
	 & & \\ 
 

    \theutterance \stepcounter{utterance}  

    & & & \multicolumn{2}{p{0.3\linewidth}}{\cellcolor[rgb]{0.95,0.95,0.95}{%
	\makecell[{{p{\linewidth}}}]{% 
	  \tt {\tiny [GM$|$GM]}  
	 COMUNICA: Ciao! Il mio tavolo ha: P in (1,2), C in (2,3), L in (5,1). Sto cercando di allineare tutti gli oggetti alla stessa posizione del tuo tavolo. Per favore, fammi sapere dove hai posizionato P, C e L, o suggeriscimi una configurazione obiettivo comune (es: "Mettiamo tutti al centro: P,C,L in (4,4)"). Possiamo farlo in 3 mosse! 
	  } 
	   } 
	   } 
	 & & \\ 
 

    \theutterance \stepcounter{utterance}  

    & & \multicolumn{4}{p{0.6\linewidth}}{\cellcolor[rgb]{0.9,0.9,0.9}{%
	\makecell[{{p{\linewidth}}}]{% 
	  \tt {\tiny [GM$\rangle$B]}  
	 Sono il tuo game master, e stai giocando un gioco collaborativo con la seguente griglia come tavolo da gioco:\\ \tt \\ \tt ```\\ \tt  1234567\\ \tt ╔═══════╗\\ \tt ║◌◌◌◌C◌◌║ 1\\ \tt ║L◌◌◌◌◌◌║ 2\\ \tt ║P◌◌◌◌┌─╢ 3\\ \tt ╟──┐◌◌│◌║ 4\\ \tt ║◌◌├──┼─╢ 5\\ \tt ║◌◌│◌◌│◌║ 6\\ \tt ║◌◌│◌◌│◌║ 7\\ \tt ╚══╧══╧═╝\\ \tt \\ \tt ```\\ \tt \\ \tt * Il bordo superiore mostra le coordinate x che aumentano verso destra, e il bordo destro le coordinate y che aumentano verso il basso.\\ \tt * I seguenti oggetti sono posizionati casualmente sulla tua griglia: 'C', 'L', 'P'.\\ \tt \\ \tt \\ \tt L'altro giocatore vede una variante del tavolo da gioco, dove gli oggetti sono posizionati in diverse posizioni casuali. Non puoi vedere il tavolo dell'altro giocatore, e loro non possono vedere il tuo.\\ \tt \\ \tt \\ \tt **Obiettivo:**\\ \tt \\ \tt Entrambi i giocatori devono spostare gli oggetti sui rispettivi sfondi in modo che oggetti identici finiscano alle stesse coordinate. Devi comunicare con l'altro giocatore per concordare una configurazione obiettivo comune.\\ \tt \\ \tt \\ \tt **Regole:**\\ \tt \\ \tt * In ogni turno, puoi inviare esattamente uno dei seguenti due comandi:\\ \tt  1. `COMUNICA: <MESSAGE>`: per inviare un messaggio (tutto fino alla prossima interruzione di riga) all'altro giocatore. Lo inoltrerò al tuo partner.\\ \tt 2. `SPOSTA: <OBJECT>, (<X>, <Y>)`: per spostare un oggetto in una nuova posizione, dove <X> è la colonna e <Y> è la riga. Ti farò sapere se la tua mossa è andata a buon fine o no.\\ \tt * Se non rispetti il formato, o invii diversi comandi contemporaneamente, devo penalizzarti.\\ \tt * Se entrambi i giocatori accumulano più di 8 penalità, perdete entrambi il gioco.\\ \tt * È vitale che comunichi con l'altro giocatore riguardo al tuo stato obiettivo! L'*unico* modo in cui puoi trasmettere la tua strategia all'altro giocatore è usando il comando `COMUNICA: <MESSAGE>`!\\ \tt \\ \tt \\ \tt **Spostamento Oggetti**\\ \tt \\ \tt * Puoi spostare oggetti solo in celle entro i limiti della griglia. La cella di destinazione deve essere vuota, cioè deve contenere solo il simbolo '◌'.\\ \tt * Se provi a spostare un oggetto in un posto che non è vuoto, o provi a spostarlo fuori dalla griglia, devo penalizzarti. Hai un altro tentativo.\\ \tt * Prima di fare una mossa, controlla due volte che il posto di destinazione sia vuoto, e non contenga alcuna lettera, cornice, o linea!\\ \tt \\ \tt \\ \tt **Fine del Gioco**\\ \tt \\ \tt Se pensi di aver raggiunto l'obiettivo di allineare tutti gli oggetti, puoi chiedere all'altro giocatore di finire il gioco inviando COMUNICA: finito?. Se l'altro giocatore ti chiede di finire il gioco, e rispondi con COMUNICA: finito!, the game will end.\\ \tt \\ \tt Entrambi i giocatori vincono se il gioco finisce entro 12 round, dove un round è definito come due giocatori che inviano ciascuno un comando valido.\\ \tt \\ \tt **Punteggio:**\\ \tt \\ \tt Più vicini sono gli oggetti identici in entrambi i tavoli da gioco, più punti ottieni. Le penalità riducono i tuoi punti. Riesci a battere il record?\\ \tt \\ \tt L'altro giocatore ha iniziato il gioco inviando questo messaggio:\\ \tt \\ \tt "Ciao! Il mio tavolo ha: P in (1,2), C in (2,3), L in (5,1). Sto cercando di allineare tutti gli oggetti alla stessa posizione del tuo tavolo. Per favore, fammi sapere dove hai posizionato P, C e L, o suggeriscimi una configurazione obiettivo comune (es: "Mettiamo tutti al centro: P,C,L in (4,4)"). Possiamo farlo in 3 mosse!"\\ \tt \\ \tt Qual è il tuo primo comando? 
	  } 
	   } 
	   } 
	 & \\ 
 

    \theutterance \stepcounter{utterance}  

    & & & \multicolumn{4}{p{0.6\linewidth}}{\cellcolor[rgb]{1,0.85,0.72}{%
	\makecell[{{p{\linewidth}}}]{% 
	  \tt {\tiny [GM$\langle$B]}  
	 COMUNICA: Ciao! Ho P in (1,3), C in (1,5), L in (1,2). La tua configurazione è molto diversa dalla mia. Suggerisco di spostare tutti gli oggetti in (4,4) — è il centro della griglia, accessibile da tutte le posizioni e facile da verificare. Possiamo farlo in 3 mosse ciascuno: io muovo P → (4,4), poi C → (4,4), poi L → (4,4). Tu fai lo stesso. Fammi sapere se sei d'accordo! 
	  } 
	   } 
	   } 
	  \\ 
 

    \theutterance \stepcounter{utterance}  

    & & & \multicolumn{2}{p{0.3\linewidth}}{\cellcolor[rgb]{0.95,0.95,0.95}{%
	\makecell[{{p{\linewidth}}}]{% 
	  \tt {\tiny [GM$|$GM]}  
	 COMUNICA: Ciao! Ho P in (1,3), C in (1,5), L in (1,2). La tua configurazione è molto diversa dalla mia. Suggerisco di spostare tutti gli oggetti in (4,4) — è il centro della griglia, accessibile da tutte le posizioni e facile da verificare. Possiamo farlo in 3 mosse ciascuno: io muovo P → (4,4), poi C → (4,4), poi L → (4,4). Tu fai lo stesso. Fammi sapere se sei d'accordo! 
	  } 
	   } 
	   } 
	 & & \\ 
 

    \theutterance \stepcounter{utterance}  

    & & \multicolumn{4}{p{0.6\linewidth}}{\cellcolor[rgb]{0.9,0.9,0.9}{%
	\makecell[{{p{\linewidth}}}]{% 
	  \tt {\tiny [A$\langle$GM]}  
	 Il tuo messaggio è stato inoltrato all'altro giocatore.\\ \tt \\ \tt \\ \tt Stai attualmente giocando il round 1 di massimo 12.\\ \tt Avete accumulato collettivamente 0 di 8 penalità.\\ \tt L'altro giocatore ha inviato questo messaggio:\\ \tt \\ \tt \\ \tt "Ciao! Ho P in (1,3), C in (1,5), L in (1,2). La tua configurazione è molto diversa dalla mia. Suggerisco di spostare tutti gli oggetti in (4,4) — è il centro della griglia, accessibile da tutte le posizioni e facile da verificare. Possiamo farlo in 3 mosse ciascuno: io muovo P → (4,4), poi C → (4,4), poi L → (4,4). Tu fai lo stesso. Fammi sapere se sei d'accordo!"\\ \tt \\ \tt \\ \tt Qual è il tuo prossimo comando? 
	  } 
	   } 
	   } 
	 & \\ 
 

    \theutterance \stepcounter{utterance}  

    & \multicolumn{4}{p{0.6\linewidth}}{\cellcolor[rgb]{0.8,1,0.9}{%
	\makecell[{{p{\linewidth}}}]{% 
	  \tt {\tiny [A$\rangle$GM]}  
	 COMUNICA: Accordo totale! Muoverò P da (1,2) → (4,4), poi C da (2,3) → (4,4), poi L da (5,1) → (4,4). Inizio con P: SPOSTA: P, (4,4) 
	  } 
	   } 
	   } 
	 & & \\ 
 

    \theutterance \stepcounter{utterance}  

    & & & \multicolumn{2}{p{0.3\linewidth}}{\cellcolor[rgb]{0.95,0.95,0.95}{%
	\makecell[{{p{\linewidth}}}]{% 
	  \tt {\tiny [GM$|$GM]}  
	 COMUNICA: Accordo totale! Muoverò P da (1,2) → (4,4), poi C da (2,3) → (4,4), poi L da (5,1) → (4,4). Inizio con P: SPOSTA: P, (4,4) 
	  } 
	   } 
	   } 
	 & & \\ 
 

    \theutterance \stepcounter{utterance}  

    & & & \multicolumn{2}{p{0.3\linewidth}}{\cellcolor[rgb]{0.95,0.95,0.95}{%
	\makecell[{{p{\linewidth}}}]{% 
	  \tt {\tiny [GM$|$GM]}  
	 Invalid response format: COMUNICA: Accordo totale! Muoverò P da (1,2) → (4,4), poi C da (2,3) → (4,4), poi L da (5,1) → (4,4). Inizio con P: SPOSTA: P, (4,4) 
	  } 
	   } 
	   } 
	 & & \\ 
 

    \theutterance \stepcounter{utterance}  

    & & \multicolumn{4}{p{0.6\linewidth}}{\cellcolor[rgb]{0.9,0.9,0.9}{%
	\makecell[{{p{\linewidth}}}]{% 
	  \tt {\tiny [A$\langle$GM]}  
	 Penalità: Il tuo messaggio contiene più di un comando!\\ \tt \\ \tt Assicurati che la tua risposta contenga solo `COMUNICA: <MESSAGE>` o `SPOSTA: <OBJECT>, (<X>, <Y>)`, e nient'altro!\\ \tt \\ \tt Avete accumulato collettivamente 1 di 8 penalità.\\ \tt Riprova! 
	  } 
	   } 
	   } 
	 & \\ 
 

    \theutterance \stepcounter{utterance}  

    & \multicolumn{4}{p{0.6\linewidth}}{\cellcolor[rgb]{0.8,1,0.9}{%
	\makecell[{{p{\linewidth}}}]{% 
	  \tt {\tiny [A$\rangle$GM]}  
	 COMUNICA: Accordo totale! Muoverò P da (1,2) → (4,4), poi C da (2,3) → (4,4), poi L da (5,1) → (4,4). Inizio con P. 
	  } 
	   } 
	   } 
	 & & \\ 
 

    \theutterance \stepcounter{utterance}  

    & & & \multicolumn{2}{p{0.3\linewidth}}{\cellcolor[rgb]{0.95,0.95,0.95}{%
	\makecell[{{p{\linewidth}}}]{% 
	  \tt {\tiny [GM$|$GM]}  
	 COMUNICA: Accordo totale! Muoverò P da (1,2) → (4,4), poi C da (2,3) → (4,4), poi L da (5,1) → (4,4). Inizio con P. 
	  } 
	   } 
	   } 
	 & & \\ 
 

    \theutterance \stepcounter{utterance}  

    & & \multicolumn{4}{p{0.6\linewidth}}{\cellcolor[rgb]{0.9,0.9,0.9}{%
	\makecell[{{p{\linewidth}}}]{% 
	  \tt {\tiny [GM$\rangle$B]}  
	 Il tuo messaggio è stato inoltrato all'altro giocatore.\\ \tt \\ \tt \\ \tt Stai attualmente giocando il round 2 di massimo 12.\\ \tt Avete accumulato collettivamente 1 di 8 penalità.\\ \tt L'altro giocatore ha inviato questo messaggio:\\ \tt \\ \tt \\ \tt "Accordo totale! Muoverò P da (1,2) → (4,4), poi C da (2,3) → (4,4), poi L da (5,1) → (4,4). Inizio con P."\\ \tt \\ \tt \\ \tt Qual è il tuo prossimo comando? 
	  } 
	   } 
	   } 
	 & \\ 
 

    \theutterance \stepcounter{utterance}  

    & & & \multicolumn{4}{p{0.6\linewidth}}{\cellcolor[rgb]{1,0.85,0.72}{%
	\makecell[{{p{\linewidth}}}]{% 
	  \tt {\tiny [GM$\langle$B]}  
	 SPOSTA: P, (4,4) 
	  } 
	   } 
	   } 
	  \\ 
 

    \theutterance \stepcounter{utterance}  

    & & & \multicolumn{2}{p{0.3\linewidth}}{\cellcolor[rgb]{0.95,0.95,0.95}{%
	\makecell[{{p{\linewidth}}}]{% 
	  \tt {\tiny [GM$|$GM]}  
	 SPOSTA: P, (4,4) 
	  } 
	   } 
	   } 
	 & & \\ 
 

    \theutterance \stepcounter{utterance}  

    & & & \multicolumn{2}{p{0.3\linewidth}}{\cellcolor[rgb]{0.95,0.95,0.95}{%
	\makecell[{{p{\linewidth}}}]{% 
	  \tt {\tiny [GM$|$GM]}  
	 Spostato 'P' in (4,4) con successo. La tua griglia aggiornata appare così:\\ \tt ```\\ \tt  1234567\\ \tt ╔═══════╗\\ \tt ║◌◌◌◌C◌◌║ 1\\ \tt ║L◌◌◌◌◌◌║ 2\\ \tt ║◌◌◌◌◌┌─╢ 3\\ \tt ╟──┐P◌│◌║ 4\\ \tt ║◌◌├──┼─╢ 5\\ \tt ║◌◌│◌◌│◌║ 6\\ \tt ║◌◌│◌◌│◌║ 7\\ \tt ╚══╧══╧═╝\\ \tt ```\\ \tt  
	  } 
	   } 
	   } 
	 & & \\ 
 

    \theutterance \stepcounter{utterance}  

    & & \multicolumn{4}{p{0.6\linewidth}}{\cellcolor[rgb]{0.9,0.9,0.9}{%
	\makecell[{{p{\linewidth}}}]{% 
	  \tt {\tiny [A$\langle$GM]}  
	 Il tuo messaggio è stato inoltrato all'altro giocatore.\\ \tt \\ \tt \\ \tt Stai attualmente giocando il round 2 di massimo 12.\\ \tt Avete accumulato collettivamente 1 di 8 penalità.\\ \tt L'altro giocatore ha spostato un oggetto sulla sua griglia.\\ \tt \\ \tt \\ \tt Qual è il tuo prossimo comando? 
	  } 
	   } 
	   } 
	 & \\ 
 

    \theutterance \stepcounter{utterance}  

    & \multicolumn{4}{p{0.6\linewidth}}{\cellcolor[rgb]{0.8,1,0.9}{%
	\makecell[{{p{\linewidth}}}]{% 
	  \tt {\tiny [A$\rangle$GM]}  
	 SPOSTA: P, (4,4) 
	  } 
	   } 
	   } 
	 & & \\ 
 

    \theutterance \stepcounter{utterance}  

    & & & \multicolumn{2}{p{0.3\linewidth}}{\cellcolor[rgb]{0.95,0.95,0.95}{%
	\makecell[{{p{\linewidth}}}]{% 
	  \tt {\tiny [GM$|$GM]}  
	 SPOSTA: P, (4,4) 
	  } 
	   } 
	   } 
	 & & \\ 
 

    \theutterance \stepcounter{utterance}  

    & & & \multicolumn{2}{p{0.3\linewidth}}{\cellcolor[rgb]{0.95,0.95,0.95}{%
	\makecell[{{p{\linewidth}}}]{% 
	  \tt {\tiny [GM$|$GM]}  
	 Spostato 'P' in (4,4) con successo. La tua griglia aggiornata appare così:\\ \tt ```\\ \tt  1234567\\ \tt ╔═══════╗\\ \tt ║◌◌◌◌L◌◌║ 1\\ \tt ║◌◌◌◌◌◌◌║ 2\\ \tt ║◌C◌◌◌┌─╢ 3\\ \tt ╟──┐P◌│◌║ 4\\ \tt ║◌◌├──┼─╢ 5\\ \tt ║◌◌│◌◌│◌║ 6\\ \tt ║◌◌│◌◌│◌║ 7\\ \tt ╚══╧══╧═╝\\ \tt ```\\ \tt  
	  } 
	   } 
	   } 
	 & & \\ 
 

    \theutterance \stepcounter{utterance}  

    & & \multicolumn{4}{p{0.6\linewidth}}{\cellcolor[rgb]{0.9,0.9,0.9}{%
	\makecell[{{p{\linewidth}}}]{% 
	  \tt {\tiny [GM$\rangle$B]}  
	 Spostato 'P' in (4,4) con successo. La tua griglia aggiornata appare così:\\ \tt ```\\ \tt  1234567\\ \tt ╔═══════╗\\ \tt ║◌◌◌◌C◌◌║ 1\\ \tt ║L◌◌◌◌◌◌║ 2\\ \tt ║◌◌◌◌◌┌─╢ 3\\ \tt ╟──┐P◌│◌║ 4\\ \tt ║◌◌├──┼─╢ 5\\ \tt ║◌◌│◌◌│◌║ 6\\ \tt ║◌◌│◌◌│◌║ 7\\ \tt ╚══╧══╧═╝\\ \tt ```\\ \tt Stai attualmente giocando il round 3 di massimo 12.\\ \tt Avete accumulato collettivamente 1 di 8 penalità.\\ \tt L'altro giocatore ha spostato un oggetto sulla sua griglia.\\ \tt \\ \tt \\ \tt Qual è il tuo prossimo comando? 
	  } 
	   } 
	   } 
	 & \\ 
 

    \theutterance \stepcounter{utterance}  

    & & & \multicolumn{4}{p{0.6\linewidth}}{\cellcolor[rgb]{1,0.85,0.72}{%
	\makecell[{{p{\linewidth}}}]{% 
	  \tt {\tiny [GM$\langle$B]}  
	 SPOSTA: C, (4,4) 
	  } 
	   } 
	   } 
	  \\ 
 

    \theutterance \stepcounter{utterance}  

    & & & \multicolumn{2}{p{0.3\linewidth}}{\cellcolor[rgb]{0.95,0.95,0.95}{%
	\makecell[{{p{\linewidth}}}]{% 
	  \tt {\tiny [GM$|$GM]}  
	 SPOSTA: C, (4,4) 
	  } 
	   } 
	   } 
	 & & \\ 
 

    \theutterance \stepcounter{utterance}  

    & & & \multicolumn{2}{p{0.3\linewidth}}{\cellcolor[rgb]{0.95,0.95,0.95}{%
	\makecell[{{p{\linewidth}}}]{% 
	  \tt {\tiny [GM$|$GM]}  
	 Penalità: (4,4) non è vuoto, ma contiene 'P'.\\ \tt \\ \tt Avete accumulato collettivamente 2 di 8 penalità.\\ \tt Riprova! 
	  } 
	   } 
	   } 
	 & & \\ 
 

    \theutterance \stepcounter{utterance}  

    & & \multicolumn{4}{p{0.6\linewidth}}{\cellcolor[rgb]{0.9,0.9,0.9}{%
	\makecell[{{p{\linewidth}}}]{% 
	  \tt {\tiny [GM$\rangle$B]}  
	 Penalità: (4,4) non è vuoto, ma contiene 'P'.\\ \tt \\ \tt Avete accumulato collettivamente 2 di 8 penalità.\\ \tt Riprova! 
	  } 
	   } 
	   } 
	 & \\ 
 

    \theutterance \stepcounter{utterance}  

    & & & \multicolumn{4}{p{0.6\linewidth}}{\cellcolor[rgb]{1,0.85,0.72}{%
	\makecell[{{p{\linewidth}}}]{% 
	  \tt {\tiny [GM$\langle$B]}  
	 SPOSTA: C, (5,4) 
	  } 
	   } 
	   } 
	  \\ 
 

    \theutterance \stepcounter{utterance}  

    & & & \multicolumn{2}{p{0.3\linewidth}}{\cellcolor[rgb]{0.95,0.95,0.95}{%
	\makecell[{{p{\linewidth}}}]{% 
	  \tt {\tiny [GM$|$GM]}  
	 SPOSTA: C, (5,4) 
	  } 
	   } 
	   } 
	 & & \\ 
 

    \theutterance \stepcounter{utterance}  

    & & & \multicolumn{2}{p{0.3\linewidth}}{\cellcolor[rgb]{0.95,0.95,0.95}{%
	\makecell[{{p{\linewidth}}}]{% 
	  \tt {\tiny [GM$|$GM]}  
	 Spostato 'C' in (5,4) con successo. La tua griglia aggiornata appare così:\\ \tt ```\\ \tt  1234567\\ \tt ╔═══════╗\\ \tt ║◌◌◌◌◌◌◌║ 1\\ \tt ║L◌◌◌◌◌◌║ 2\\ \tt ║◌◌◌◌◌┌─╢ 3\\ \tt ╟──┐PC│◌║ 4\\ \tt ║◌◌├──┼─╢ 5\\ \tt ║◌◌│◌◌│◌║ 6\\ \tt ║◌◌│◌◌│◌║ 7\\ \tt ╚══╧══╧═╝\\ \tt ```\\ \tt  
	  } 
	   } 
	   } 
	 & & \\ 
 

    \theutterance \stepcounter{utterance}  

    & & \multicolumn{4}{p{0.6\linewidth}}{\cellcolor[rgb]{0.9,0.9,0.9}{%
	\makecell[{{p{\linewidth}}}]{% 
	  \tt {\tiny [A$\langle$GM]}  
	 Spostato 'P' in (4,4) con successo. La tua griglia aggiornata appare così:\\ \tt ```\\ \tt  1234567\\ \tt ╔═══════╗\\ \tt ║◌◌◌◌L◌◌║ 1\\ \tt ║◌◌◌◌◌◌◌║ 2\\ \tt ║◌C◌◌◌┌─╢ 3\\ \tt ╟──┐P◌│◌║ 4\\ \tt ║◌◌├──┼─╢ 5\\ \tt ║◌◌│◌◌│◌║ 6\\ \tt ║◌◌│◌◌│◌║ 7\\ \tt ╚══╧══╧═╝\\ \tt ```\\ \tt Stai attualmente giocando il round 3 di massimo 12.\\ \tt Avete accumulato collettivamente 2 di 8 penalità.\\ \tt L'altro giocatore ha spostato un oggetto sulla sua griglia.\\ \tt \\ \tt \\ \tt Qual è il tuo prossimo comando? 
	  } 
	   } 
	   } 
	 & \\ 
 

    \theutterance \stepcounter{utterance}  

    & \multicolumn{4}{p{0.6\linewidth}}{\cellcolor[rgb]{0.8,1,0.9}{%
	\makecell[{{p{\linewidth}}}]{% 
	  \tt {\tiny [A$\rangle$GM]}  
	 COMUNICA: Ho spostato P in (4,4). Ora muovo C da (2,3) → (4,4). Aspetto il tuo movimento su C e L. Quando tutti e tre sono in (4,4), avremo vinto! 
	  } 
	   } 
	   } 
	 & & \\ 
 

    \theutterance \stepcounter{utterance}  

    & & & \multicolumn{2}{p{0.3\linewidth}}{\cellcolor[rgb]{0.95,0.95,0.95}{%
	\makecell[{{p{\linewidth}}}]{% 
	  \tt {\tiny [GM$|$GM]}  
	 COMUNICA: Ho spostato P in (4,4). Ora muovo C da (2,3) → (4,4). Aspetto il tuo movimento su C e L. Quando tutti e tre sono in (4,4), avremo vinto! 
	  } 
	   } 
	   } 
	 & & \\ 
 

    \theutterance \stepcounter{utterance}  

    & & \multicolumn{4}{p{0.6\linewidth}}{\cellcolor[rgb]{0.9,0.9,0.9}{%
	\makecell[{{p{\linewidth}}}]{% 
	  \tt {\tiny [GM$\rangle$B]}  
	 Spostato 'C' in (5,4) con successo. La tua griglia aggiornata appare così:\\ \tt ```\\ \tt  1234567\\ \tt ╔═══════╗\\ \tt ║◌◌◌◌◌◌◌║ 1\\ \tt ║L◌◌◌◌◌◌║ 2\\ \tt ║◌◌◌◌◌┌─╢ 3\\ \tt ╟──┐PC│◌║ 4\\ \tt ║◌◌├──┼─╢ 5\\ \tt ║◌◌│◌◌│◌║ 6\\ \tt ║◌◌│◌◌│◌║ 7\\ \tt ╚══╧══╧═╝\\ \tt ```\\ \tt Stai attualmente giocando il round 4 di massimo 12.\\ \tt Avete accumulato collettivamente 2 di 8 penalità.\\ \tt L'altro giocatore ha inviato questo messaggio:\\ \tt \\ \tt \\ \tt "Ho spostato P in (4,4). Ora muovo C da (2,3) → (4,4). Aspetto il tuo movimento su C e L. Quando tutti e tre sono in (4,4), avremo vinto!"\\ \tt \\ \tt \\ \tt Qual è il tuo prossimo comando? 
	  } 
	   } 
	   } 
	 & \\ 
 

    \theutterance \stepcounter{utterance}  

    & & & \multicolumn{4}{p{0.6\linewidth}}{\cellcolor[rgb]{1,0.85,0.72}{%
	\makecell[{{p{\linewidth}}}]{% 
	  \tt {\tiny [GM$\langle$B]}  
	 SPOSTA: L, (4,4) 
	  } 
	   } 
	   } 
	  \\ 
 

    \theutterance \stepcounter{utterance}  

    & & & \multicolumn{2}{p{0.3\linewidth}}{\cellcolor[rgb]{0.95,0.95,0.95}{%
	\makecell[{{p{\linewidth}}}]{% 
	  \tt {\tiny [GM$|$GM]}  
	 SPOSTA: L, (4,4) 
	  } 
	   } 
	   } 
	 & & \\ 
 

    \theutterance \stepcounter{utterance}  

    & & & \multicolumn{2}{p{0.3\linewidth}}{\cellcolor[rgb]{0.95,0.95,0.95}{%
	\makecell[{{p{\linewidth}}}]{% 
	  \tt {\tiny [GM$|$GM]}  
	 Penalità: (4,4) non è vuoto, ma contiene 'P'.\\ \tt \\ \tt Avete accumulato collettivamente 3 di 8 penalità.\\ \tt Riprova! 
	  } 
	   } 
	   } 
	 & & \\ 
 

    \theutterance \stepcounter{utterance}  

    & & \multicolumn{4}{p{0.6\linewidth}}{\cellcolor[rgb]{0.9,0.9,0.9}{%
	\makecell[{{p{\linewidth}}}]{% 
	  \tt {\tiny [GM$\rangle$B]}  
	 Penalità: (4,4) non è vuoto, ma contiene 'P'.\\ \tt \\ \tt Avete accumulato collettivamente 3 di 8 penalità.\\ \tt Riprova! 
	  } 
	   } 
	   } 
	 & \\ 
 

    \theutterance \stepcounter{utterance}  

    & & & \multicolumn{4}{p{0.6\linewidth}}{\cellcolor[rgb]{1,0.85,0.72}{%
	\makecell[{{p{\linewidth}}}]{% 
	  \tt {\tiny [GM$\langle$B]}  
	 SPOSTA: L, (3,4) 
	  } 
	   } 
	   } 
	  \\ 
 

    \theutterance \stepcounter{utterance}  

    & & & \multicolumn{2}{p{0.3\linewidth}}{\cellcolor[rgb]{0.95,0.95,0.95}{%
	\makecell[{{p{\linewidth}}}]{% 
	  \tt {\tiny [GM$|$GM]}  
	 SPOSTA: L, (3,4) 
	  } 
	   } 
	   } 
	 & & \\ 
 

    \theutterance \stepcounter{utterance}  

    & & & \multicolumn{2}{p{0.3\linewidth}}{\cellcolor[rgb]{0.95,0.95,0.95}{%
	\makecell[{{p{\linewidth}}}]{% 
	  \tt {\tiny [GM$|$GM]}  
	 Penalità: (3,4) non è vuoto, ma contiene '┐'.\\ \tt \\ \tt Avete accumulato collettivamente 4 di 8 penalità.\\ \tt Riprova! 
	  } 
	   } 
	   } 
	 & & \\ 
 

    \theutterance \stepcounter{utterance}  

    & & \multicolumn{4}{p{0.6\linewidth}}{\cellcolor[rgb]{0.9,0.9,0.9}{%
	\makecell[{{p{\linewidth}}}]{% 
	  \tt {\tiny [GM$\rangle$B]}  
	 Penalità: (3,4) non è vuoto, ma contiene '┐'.\\ \tt \\ \tt Avete accumulato collettivamente 4 di 8 penalità.\\ \tt Riprova! 
	  } 
	   } 
	   } 
	 & \\ 
 

    \theutterance \stepcounter{utterance}  

    & & & \multicolumn{4}{p{0.6\linewidth}}{\cellcolor[rgb]{1,0.85,0.72}{%
	\makecell[{{p{\linewidth}}}]{% 
	  \tt {\tiny [GM$\langle$B]}  
	 SPOSTA: L, (4,3) 
	  } 
	   } 
	   } 
	  \\ 
 

    \theutterance \stepcounter{utterance}  

    & & & \multicolumn{2}{p{0.3\linewidth}}{\cellcolor[rgb]{0.95,0.95,0.95}{%
	\makecell[{{p{\linewidth}}}]{% 
	  \tt {\tiny [GM$|$GM]}  
	 SPOSTA: L, (4,3) 
	  } 
	   } 
	   } 
	 & & \\ 
 

    \theutterance \stepcounter{utterance}  

    & & & \multicolumn{2}{p{0.3\linewidth}}{\cellcolor[rgb]{0.95,0.95,0.95}{%
	\makecell[{{p{\linewidth}}}]{% 
	  \tt {\tiny [GM$|$GM]}  
	 Spostato 'L' in (4,3) con successo. La tua griglia aggiornata appare così:\\ \tt ```\\ \tt  1234567\\ \tt ╔═══════╗\\ \tt ║◌◌◌◌◌◌◌║ 1\\ \tt ║◌◌◌◌◌◌◌║ 2\\ \tt ║◌◌◌L◌┌─╢ 3\\ \tt ╟──┐PC│◌║ 4\\ \tt ║◌◌├──┼─╢ 5\\ \tt ║◌◌│◌◌│◌║ 6\\ \tt ║◌◌│◌◌│◌║ 7\\ \tt ╚══╧══╧═╝\\ \tt ```\\ \tt  
	  } 
	   } 
	   } 
	 & & \\ 
 

    \theutterance \stepcounter{utterance}  

    & & \multicolumn{4}{p{0.6\linewidth}}{\cellcolor[rgb]{0.9,0.9,0.9}{%
	\makecell[{{p{\linewidth}}}]{% 
	  \tt {\tiny [A$\langle$GM]}  
	 Il tuo messaggio è stato inoltrato all'altro giocatore.\\ \tt \\ \tt \\ \tt Stai attualmente giocando il round 4 di massimo 12.\\ \tt Avete accumulato collettivamente 4 di 8 penalità.\\ \tt L'altro giocatore ha spostato un oggetto sulla sua griglia.\\ \tt \\ \tt \\ \tt Qual è il tuo prossimo comando? 
	  } 
	   } 
	   } 
	 & \\ 
 

    \theutterance \stepcounter{utterance}  

    & \multicolumn{4}{p{0.6\linewidth}}{\cellcolor[rgb]{0.8,1,0.9}{%
	\makecell[{{p{\linewidth}}}]{% 
	  \tt {\tiny [A$\rangle$GM]}  
	 SPOSTA: C, (4,4) 
	  } 
	   } 
	   } 
	 & & \\ 
 

    \theutterance \stepcounter{utterance}  

    & & & \multicolumn{2}{p{0.3\linewidth}}{\cellcolor[rgb]{0.95,0.95,0.95}{%
	\makecell[{{p{\linewidth}}}]{% 
	  \tt {\tiny [GM$|$GM]}  
	 SPOSTA: C, (4,4) 
	  } 
	   } 
	   } 
	 & & \\ 
 

    \theutterance \stepcounter{utterance}  

    & & & \multicolumn{2}{p{0.3\linewidth}}{\cellcolor[rgb]{0.95,0.95,0.95}{%
	\makecell[{{p{\linewidth}}}]{% 
	  \tt {\tiny [GM$|$GM]}  
	 Penalità: (4,4) non è vuoto, ma contiene 'P'.\\ \tt \\ \tt Avete accumulato collettivamente 5 di 8 penalità.\\ \tt Riprova! 
	  } 
	   } 
	   } 
	 & & \\ 
 

    \theutterance \stepcounter{utterance}  

    & & \multicolumn{4}{p{0.6\linewidth}}{\cellcolor[rgb]{0.9,0.9,0.9}{%
	\makecell[{{p{\linewidth}}}]{% 
	  \tt {\tiny [A$\langle$GM]}  
	 Penalità: (4,4) non è vuoto, ma contiene 'P'.\\ \tt \\ \tt Avete accumulato collettivamente 5 di 8 penalità.\\ \tt Riprova! 
	  } 
	   } 
	   } 
	 & \\ 
 

    \theutterance \stepcounter{utterance}  

    & \multicolumn{4}{p{0.6\linewidth}}{\cellcolor[rgb]{0.8,1,0.9}{%
	\makecell[{{p{\linewidth}}}]{% 
	  \tt {\tiny [A$\rangle$GM]}  
	 COMUNICA: Oops! (4,4) è occupato da P. Cambiamo obiettivo: spostiamo tutti in (3,4) — è vuoto e vicino. Muoverò C a (3,4), poi L a (3,4), e ti chiederò di fare lo stesso. Confermi (3,4) come nuovo obiettivo? 
	  } 
	   } 
	   } 
	 & & \\ 
 

    \theutterance \stepcounter{utterance}  

    & & & \multicolumn{2}{p{0.3\linewidth}}{\cellcolor[rgb]{0.95,0.95,0.95}{%
	\makecell[{{p{\linewidth}}}]{% 
	  \tt {\tiny [GM$|$GM]}  
	 COMUNICA: Oops! (4,4) è occupato da P. Cambiamo obiettivo: spostiamo tutti in (3,4) — è vuoto e vicino. Muoverò C a (3,4), poi L a (3,4), e ti chiederò di fare lo stesso. Confermi (3,4) come nuovo obiettivo? 
	  } 
	   } 
	   } 
	 & & \\ 
 

    \theutterance \stepcounter{utterance}  

    & & \multicolumn{4}{p{0.6\linewidth}}{\cellcolor[rgb]{0.9,0.9,0.9}{%
	\makecell[{{p{\linewidth}}}]{% 
	  \tt {\tiny [GM$\rangle$B]}  
	 Spostato 'L' in (4,3) con successo. La tua griglia aggiornata appare così:\\ \tt ```\\ \tt  1234567\\ \tt ╔═══════╗\\ \tt ║◌◌◌◌◌◌◌║ 1\\ \tt ║◌◌◌◌◌◌◌║ 2\\ \tt ║◌◌◌L◌┌─╢ 3\\ \tt ╟──┐PC│◌║ 4\\ \tt ║◌◌├──┼─╢ 5\\ \tt ║◌◌│◌◌│◌║ 6\\ \tt ║◌◌│◌◌│◌║ 7\\ \tt ╚══╧══╧═╝\\ \tt ```\\ \tt Stai attualmente giocando il round 5 di massimo 12.\\ \tt Avete accumulato collettivamente 5 di 8 penalità.\\ \tt L'altro giocatore ha inviato questo messaggio:\\ \tt \\ \tt \\ \tt "Oops! (4,4) è occupato da P. Cambiamo obiettivo: spostiamo tutti in (3,4) — è vuoto e vicino. Muoverò C a (3,4), poi L a (3,4), e ti chiederò di fare lo stesso. Confermi (3,4) come nuovo obiettivo?"\\ \tt \\ \tt \\ \tt Qual è il tuo prossimo comando? 
	  } 
	   } 
	   } 
	 & \\ 
 

    \theutterance \stepcounter{utterance}  

    & & & \multicolumn{4}{p{0.6\linewidth}}{\cellcolor[rgb]{1,0.85,0.72}{%
	\makecell[{{p{\linewidth}}}]{% 
	  \tt {\tiny [GM$\langle$B]}  
	 COMUNICA: Confermo! Spostiamo tutti in (3,4). Ho già L in (4,3), quindi lo muovo a (3,4) ora. Attendo il tuo movimento di C verso (3,4). 
	  } 
	   } 
	   } 
	  \\ 
 

    \theutterance \stepcounter{utterance}  

    & & & \multicolumn{2}{p{0.3\linewidth}}{\cellcolor[rgb]{0.95,0.95,0.95}{%
	\makecell[{{p{\linewidth}}}]{% 
	  \tt {\tiny [GM$|$GM]}  
	 COMUNICA: Confermo! Spostiamo tutti in (3,4). Ho già L in (4,3), quindi lo muovo a (3,4) ora. Attendo il tuo movimento di C verso (3,4). 
	  } 
	   } 
	   } 
	 & & \\ 
 

    \theutterance \stepcounter{utterance}  

    & & \multicolumn{4}{p{0.6\linewidth}}{\cellcolor[rgb]{0.9,0.9,0.9}{%
	\makecell[{{p{\linewidth}}}]{% 
	  \tt {\tiny [A$\langle$GM]}  
	 Il tuo messaggio è stato inoltrato all'altro giocatore.\\ \tt \\ \tt \\ \tt Stai attualmente giocando il round 5 di massimo 12.\\ \tt Avete accumulato collettivamente 5 di 8 penalità.\\ \tt L'altro giocatore ha inviato questo messaggio:\\ \tt \\ \tt \\ \tt "Confermo! Spostiamo tutti in (3,4). Ho già L in (4,3), quindi lo muovo a (3,4) ora. Attendo il tuo movimento di C verso (3,4)."\\ \tt \\ \tt \\ \tt Qual è il tuo prossimo comando? 
	  } 
	   } 
	   } 
	 & \\ 
 

    \theutterance \stepcounter{utterance}  

    & \multicolumn{4}{p{0.6\linewidth}}{\cellcolor[rgb]{0.8,1,0.9}{%
	\makecell[{{p{\linewidth}}}]{% 
	  \tt {\tiny [A$\rangle$GM]}  
	 SPOSTA: C, (3,4) 
	  } 
	   } 
	   } 
	 & & \\ 
 

    \theutterance \stepcounter{utterance}  

    & & & \multicolumn{2}{p{0.3\linewidth}}{\cellcolor[rgb]{0.95,0.95,0.95}{%
	\makecell[{{p{\linewidth}}}]{% 
	  \tt {\tiny [GM$|$GM]}  
	 SPOSTA: C, (3,4) 
	  } 
	   } 
	   } 
	 & & \\ 
 

    \theutterance \stepcounter{utterance}  

    & & & \multicolumn{2}{p{0.3\linewidth}}{\cellcolor[rgb]{0.95,0.95,0.95}{%
	\makecell[{{p{\linewidth}}}]{% 
	  \tt {\tiny [GM$|$GM]}  
	 Penalità: (3,4) non è vuoto, ma contiene '┐'.\\ \tt \\ \tt Avete accumulato collettivamente 6 di 8 penalità.\\ \tt Riprova! 
	  } 
	   } 
	   } 
	 & & \\ 
 

    \theutterance \stepcounter{utterance}  

    & & \multicolumn{4}{p{0.6\linewidth}}{\cellcolor[rgb]{0.9,0.9,0.9}{%
	\makecell[{{p{\linewidth}}}]{% 
	  \tt {\tiny [A$\langle$GM]}  
	 Penalità: (3,4) non è vuoto, ma contiene '┐'.\\ \tt \\ \tt Avete accumulato collettivamente 6 di 8 penalità.\\ \tt Riprova! 
	  } 
	   } 
	   } 
	 & \\ 
 

    \theutterance \stepcounter{utterance}  

    & \multicolumn{4}{p{0.6\linewidth}}{\cellcolor[rgb]{0.8,1,0.9}{%
	\makecell[{{p{\linewidth}}}]{% 
	  \tt {\tiny [A$\rangle$GM]}  
	 COMUNICA: (3,4) è occupato da una parete ('┐') — non è accessibile! Propongo (4,5) come nuovo obiettivo: è vuoto, centrale, e facilmente raggiungibile da tutte le posizioni. Muoverò C verso (4,5). Se concordi, sposta anche L e P lì! 
	  } 
	   } 
	   } 
	 & & \\ 
 

    \theutterance \stepcounter{utterance}  

    & & & \multicolumn{2}{p{0.3\linewidth}}{\cellcolor[rgb]{0.95,0.95,0.95}{%
	\makecell[{{p{\linewidth}}}]{% 
	  \tt {\tiny [GM$|$GM]}  
	 COMUNICA: (3,4) è occupato da una parete ('┐') — non è accessibile! Propongo (4,5) come nuovo obiettivo: è vuoto, centrale, e facilmente raggiungibile da tutte le posizioni. Muoverò C verso (4,5). Se concordi, sposta anche L e P lì! 
	  } 
	   } 
	   } 
	 & & \\ 
 

    \theutterance \stepcounter{utterance}  

    & & \multicolumn{4}{p{0.6\linewidth}}{\cellcolor[rgb]{0.9,0.9,0.9}{%
	\makecell[{{p{\linewidth}}}]{% 
	  \tt {\tiny [GM$\rangle$B]}  
	 Il tuo messaggio è stato inoltrato all'altro giocatore.\\ \tt \\ \tt \\ \tt Stai attualmente giocando il round 6 di massimo 12.\\ \tt Avete accumulato collettivamente 6 di 8 penalità.\\ \tt L'altro giocatore ha inviato questo messaggio:\\ \tt \\ \tt \\ \tt "(3,4) è occupato da una parete ('┐') — non è accessibile! Propongo (4,5) come nuovo obiettivo: è vuoto, centrale, e facilmente raggiungibile da tutte le posizioni. Muoverò C verso (4,5). Se concordi, sposta anche L e P lì!"\\ \tt \\ \tt \\ \tt Qual è il tuo prossimo comando? 
	  } 
	   } 
	   } 
	 & \\ 
 

    \theutterance \stepcounter{utterance}  

    & & & \multicolumn{4}{p{0.6\linewidth}}{\cellcolor[rgb]{1,0.85,0.72}{%
	\makecell[{{p{\linewidth}}}]{% 
	  \tt {\tiny [GM$\langle$B]}  
	 COMUNICA: Accettato! Spostiamo tutti in (4,5). È perfetto: vuoto, centrale e accessibile. Muovo subito L da (4,3) → (4,5). Poi mi occuperò di P. 
	  } 
	   } 
	   } 
	  \\ 
 

    \theutterance \stepcounter{utterance}  

    & & & \multicolumn{2}{p{0.3\linewidth}}{\cellcolor[rgb]{0.95,0.95,0.95}{%
	\makecell[{{p{\linewidth}}}]{% 
	  \tt {\tiny [GM$|$GM]}  
	 COMUNICA: Accettato! Spostiamo tutti in (4,5). È perfetto: vuoto, centrale e accessibile. Muovo subito L da (4,3) → (4,5). Poi mi occuperò di P. 
	  } 
	   } 
	   } 
	 & & \\ 
 

    \theutterance \stepcounter{utterance}  

    & & \multicolumn{4}{p{0.6\linewidth}}{\cellcolor[rgb]{0.9,0.9,0.9}{%
	\makecell[{{p{\linewidth}}}]{% 
	  \tt {\tiny [A$\langle$GM]}  
	 Il tuo messaggio è stato inoltrato all'altro giocatore.\\ \tt \\ \tt \\ \tt Stai attualmente giocando il round 6 di massimo 12.\\ \tt Avete accumulato collettivamente 6 di 8 penalità.\\ \tt L'altro giocatore ha inviato questo messaggio:\\ \tt \\ \tt \\ \tt "Accettato! Spostiamo tutti in (4,5). È perfetto: vuoto, centrale e accessibile. Muovo subito L da (4,3) → (4,5). Poi mi occuperò di P."\\ \tt \\ \tt \\ \tt Qual è il tuo prossimo comando? 
	  } 
	   } 
	   } 
	 & \\ 
 

    \theutterance \stepcounter{utterance}  

    & \multicolumn{4}{p{0.6\linewidth}}{\cellcolor[rgb]{0.8,1,0.9}{%
	\makecell[{{p{\linewidth}}}]{% 
	  \tt {\tiny [A$\rangle$GM]}  
	 SPOSTA: C, (4,5) 
	  } 
	   } 
	   } 
	 & & \\ 
 

    \theutterance \stepcounter{utterance}  

    & & & \multicolumn{2}{p{0.3\linewidth}}{\cellcolor[rgb]{0.95,0.95,0.95}{%
	\makecell[{{p{\linewidth}}}]{% 
	  \tt {\tiny [GM$|$GM]}  
	 SPOSTA: C, (4,5) 
	  } 
	   } 
	   } 
	 & & \\ 
 

    \theutterance \stepcounter{utterance}  

    & & & \multicolumn{2}{p{0.3\linewidth}}{\cellcolor[rgb]{0.95,0.95,0.95}{%
	\makecell[{{p{\linewidth}}}]{% 
	  \tt {\tiny [GM$|$GM]}  
	 Penalità: (4,5) non è vuoto, ma contiene '─'.\\ \tt \\ \tt Avete accumulato collettivamente 7 di 8 penalità.\\ \tt Riprova! 
	  } 
	   } 
	   } 
	 & & \\ 
 

    \theutterance \stepcounter{utterance}  

    & & \multicolumn{4}{p{0.6\linewidth}}{\cellcolor[rgb]{0.9,0.9,0.9}{%
	\makecell[{{p{\linewidth}}}]{% 
	  \tt {\tiny [A$\langle$GM]}  
	 Penalità: (4,5) non è vuoto, ma contiene '─'.\\ \tt \\ \tt Avete accumulato collettivamente 7 di 8 penalità.\\ \tt Riprova! 
	  } 
	   } 
	   } 
	 & \\ 
 

    \theutterance \stepcounter{utterance}  

    & \multicolumn{4}{p{0.6\linewidth}}{\cellcolor[rgb]{0.8,1,0.9}{%
	\makecell[{{p{\linewidth}}}]{% 
	  \tt {\tiny [A$\rangle$GM]}  
	 COMUNICA: (4,5) è occupato da una linea orizzontale ('─') — non è vuoto! Propongo (5,4) come obiettivo finale: è vuoto ('◌'), centralmente posizionato, e raggiungibile da tutte le posizioni. Muoverò C verso (5,4). Confermi (5,4)? Se sì, ti chiedo di muovere P e L lì! 
	  } 
	   } 
	   } 
	 & & \\ 
 

    \theutterance \stepcounter{utterance}  

    & & & \multicolumn{2}{p{0.3\linewidth}}{\cellcolor[rgb]{0.95,0.95,0.95}{%
	\makecell[{{p{\linewidth}}}]{% 
	  \tt {\tiny [GM$|$GM]}  
	 COMUNICA: (4,5) è occupato da una linea orizzontale ('─') — non è vuoto! Propongo (5,4) come obiettivo finale: è vuoto ('◌'), centralmente posizionato, e raggiungibile da tutte le posizioni. Muoverò C verso (5,4). Confermi (5,4)? Se sì, ti chiedo di muovere P e L lì! 
	  } 
	   } 
	   } 
	 & & \\ 
 

    \theutterance \stepcounter{utterance}  

    & & \multicolumn{4}{p{0.6\linewidth}}{\cellcolor[rgb]{0.9,0.9,0.9}{%
	\makecell[{{p{\linewidth}}}]{% 
	  \tt {\tiny [GM$\rangle$B]}  
	 Il tuo messaggio è stato inoltrato all'altro giocatore.\\ \tt \\ \tt \\ \tt Stai attualmente giocando il round 7 di massimo 12.\\ \tt Avete accumulato collettivamente 7 di 8 penalità.\\ \tt L'altro giocatore ha inviato questo messaggio:\\ \tt \\ \tt \\ \tt "(4,5) è occupato da una linea orizzontale ('─') — non è vuoto! Propongo (5,4) come obiettivo finale: è vuoto ('◌'), centralmente posizionato, e raggiungibile da tutte le posizioni. Muoverò C verso (5,4). Confermi (5,4)? Se sì, ti chiedo di muovere P e L lì!"\\ \tt \\ \tt \\ \tt Qual è il tuo prossimo comando? 
	  } 
	   } 
	   } 
	 & \\ 
 

    \theutterance \stepcounter{utterance}  

    & & & \multicolumn{4}{p{0.6\linewidth}}{\cellcolor[rgb]{1,0.85,0.72}{%
	\makecell[{{p{\linewidth}}}]{% 
	  \tt {\tiny [GM$\langle$B]}  
	 COMUNICA: Confermo assolutamente (5,4)! È vuoto, perfetto e strategico. Muovo subito L da (4,3) → (5,4), poi P da (4,4) → (5,4), e finalmente C è già in (5,4) sulla tua griglia — quando lo sposti, saremo allineati! 
	  } 
	   } 
	   } 
	  \\ 
 

    \theutterance \stepcounter{utterance}  

    & & & \multicolumn{2}{p{0.3\linewidth}}{\cellcolor[rgb]{0.95,0.95,0.95}{%
	\makecell[{{p{\linewidth}}}]{% 
	  \tt {\tiny [GM$|$GM]}  
	 COMUNICA: Confermo assolutamente (5,4)! È vuoto, perfetto e strategico. Muovo subito L da (4,3) → (5,4), poi P da (4,4) → (5,4), e finalmente C è già in (5,4) sulla tua griglia — quando lo sposti, saremo allineati! 
	  } 
	   } 
	   } 
	 & & \\ 
 

    \theutterance \stepcounter{utterance}  

    & & \multicolumn{4}{p{0.6\linewidth}}{\cellcolor[rgb]{0.9,0.9,0.9}{%
	\makecell[{{p{\linewidth}}}]{% 
	  \tt {\tiny [A$\langle$GM]}  
	 Il tuo messaggio è stato inoltrato all'altro giocatore.\\ \tt \\ \tt \\ \tt Stai attualmente giocando il round 7 di massimo 12.\\ \tt Avete accumulato collettivamente 7 di 8 penalità.\\ \tt L'altro giocatore ha inviato questo messaggio:\\ \tt \\ \tt \\ \tt "Confermo assolutamente (5,4)! È vuoto, perfetto e strategico. Muovo subito L da (4,3) → (5,4), poi P da (4,4) → (5,4), e finalmente C è già in (5,4) sulla tua griglia — quando lo sposti, saremo allineati!"\\ \tt \\ \tt \\ \tt Qual è il tuo prossimo comando? 
	  } 
	   } 
	   } 
	 & \\ 
 

    \theutterance \stepcounter{utterance}  

    & \multicolumn{4}{p{0.6\linewidth}}{\cellcolor[rgb]{0.8,1,0.9}{%
	\makecell[{{p{\linewidth}}}]{% 
	  \tt {\tiny [A$\rangle$GM]}  
	 SPOSTA: C, (5,4) 
	  } 
	   } 
	   } 
	 & & \\ 
 

    \theutterance \stepcounter{utterance}  

    & & & \multicolumn{2}{p{0.3\linewidth}}{\cellcolor[rgb]{0.95,0.95,0.95}{%
	\makecell[{{p{\linewidth}}}]{% 
	  \tt {\tiny [GM$|$GM]}  
	 SPOSTA: C, (5,4) 
	  } 
	   } 
	   } 
	 & & \\ 
 

    \theutterance \stepcounter{utterance}  

    & & & \multicolumn{2}{p{0.3\linewidth}}{\cellcolor[rgb]{0.95,0.95,0.95}{%
	\makecell[{{p{\linewidth}}}]{% 
	  \tt {\tiny [GM$|$GM]}  
	 Spostato 'C' in (5,4) con successo. La tua griglia aggiornata appare così:\\ \tt ```\\ \tt  1234567\\ \tt ╔═══════╗\\ \tt ║◌◌◌◌L◌◌║ 1\\ \tt ║◌◌◌◌◌◌◌║ 2\\ \tt ║◌◌◌◌◌┌─╢ 3\\ \tt ╟──┐PC│◌║ 4\\ \tt ║◌◌├──┼─╢ 5\\ \tt ║◌◌│◌◌│◌║ 6\\ \tt ║◌◌│◌◌│◌║ 7\\ \tt ╚══╧══╧═╝\\ \tt ```\\ \tt  
	  } 
	   } 
	   } 
	 & & \\ 
 

    \theutterance \stepcounter{utterance}  

    & & \multicolumn{4}{p{0.6\linewidth}}{\cellcolor[rgb]{0.9,0.9,0.9}{%
	\makecell[{{p{\linewidth}}}]{% 
	  \tt {\tiny [GM$\rangle$B]}  
	 Il tuo messaggio è stato inoltrato all'altro giocatore.\\ \tt \\ \tt \\ \tt Stai attualmente giocando il round 8 di massimo 12.\\ \tt Avete accumulato collettivamente 7 di 8 penalità.\\ \tt L'altro giocatore ha spostato un oggetto sulla sua griglia.\\ \tt \\ \tt \\ \tt Qual è il tuo prossimo comando? 
	  } 
	   } 
	   } 
	 & \\ 
 

    \theutterance \stepcounter{utterance}  

    & & & \multicolumn{4}{p{0.6\linewidth}}{\cellcolor[rgb]{1,0.85,0.72}{%
	\makecell[{{p{\linewidth}}}]{% 
	  \tt {\tiny [GM$\langle$B]}  
	 SPOSTA: L, (5,4) 
	  } 
	   } 
	   } 
	  \\ 
 

    \theutterance \stepcounter{utterance}  

    & & & \multicolumn{2}{p{0.3\linewidth}}{\cellcolor[rgb]{0.95,0.95,0.95}{%
	\makecell[{{p{\linewidth}}}]{% 
	  \tt {\tiny [GM$|$GM]}  
	 SPOSTA: L, (5,4) 
	  } 
	   } 
	   } 
	 & & \\ 
 

    \theutterance \stepcounter{utterance}  

    & & & \multicolumn{2}{p{0.3\linewidth}}{\cellcolor[rgb]{0.95,0.95,0.95}{%
	\makecell[{{p{\linewidth}}}]{% 
	  \tt {\tiny [GM$|$GM]}  
	 Penalità: (5,4) non è vuoto, ma contiene 'C'.\\ \tt \\ \tt Avete accumulato collettivamente 8 di 8 penalità.\\ \tt Riprova! 
	  } 
	   } 
	   } 
	 & & \\ 
 

    \theutterance \stepcounter{utterance}  

    & & \multicolumn{4}{p{0.6\linewidth}}{\cellcolor[rgb]{0.9,0.9,0.9}{%
	\makecell[{{p{\linewidth}}}]{% 
	  \tt {\tiny [GM$\rangle$B]}  
	 Penalità: (5,4) non è vuoto, ma contiene 'C'.\\ \tt \\ \tt Avete accumulato collettivamente 8 di 8 penalità.\\ \tt Riprova! 
	  } 
	   } 
	   } 
	 & \\ 
 

    \theutterance \stepcounter{utterance}  

    & & & \multicolumn{4}{p{0.6\linewidth}}{\cellcolor[rgb]{1,0.85,0.72}{%
	\makecell[{{p{\linewidth}}}]{% 
	  \tt {\tiny [GM$\langle$B]}  
	 SPOSTA: P, (5,4) 
	  } 
	   } 
	   } 
	  \\ 
 

    \theutterance \stepcounter{utterance}  

    & & & \multicolumn{2}{p{0.3\linewidth}}{\cellcolor[rgb]{0.95,0.95,0.95}{%
	\makecell[{{p{\linewidth}}}]{% 
	  \tt {\tiny [GM$|$GM]}  
	 SPOSTA: P, (5,4) 
	  } 
	   } 
	   } 
	 & & \\ 
 

    \theutterance \stepcounter{utterance}  

    & & & \multicolumn{2}{p{0.3\linewidth}}{\cellcolor[rgb]{0.95,0.95,0.95}{%
	\makecell[{{p{\linewidth}}}]{% 
	  \tt {\tiny [GM$|$GM]}  
	 Penalità: (5,4) non è vuoto, ma contiene 'C'.\\ \tt \\ \tt Avete accumulato collettivamente 9 di 8 penalità.\\ \tt Riprova! 
	  } 
	   } 
	   } 
	 & & \\ 
 

    \theutterance \stepcounter{utterance}  

    & & & \multicolumn{2}{p{0.3\linewidth}}{\cellcolor[rgb]{0.95,0.95,0.95}{%
	\makecell[{{p{\linewidth}}}]{% 
	  \tt {\tiny [GM$|$GM]}  
	 Maximum number of penalties exceeded 
	  } 
	   } 
	   } 
	 & & \\ 
 

    \theutterance \stepcounter{utterance}  

    & & & \multicolumn{2}{p{0.3\linewidth}}{\cellcolor[rgb]{0.95,0.95,0.95}{%
	\makecell[{{p{\linewidth}}}]{% 
	  \tt {\tiny [GM$|$GM]}  
	 * success: False\\ \tt * lose: True\\ \tt * aborted: True\\ \tt -------\\ \tt * Shifts: 3.00\\ \tt * Max Shifts: 4.00\\ \tt * Min Shifts: 2.00\\ \tt * End Distance Sum: 2.24\\ \tt * Init Distance Sum: 8.73\\ \tt * Expected Distance Sum: 12.57\\ \tt * Penalties: 9.00\\ \tt * Max Penalties: 8.00\\ \tt * Rounds: 7.00\\ \tt * Max Rounds: 12.00\\ \tt * Object Count: 3.00\\ \tt  
	  } 
	   } 
	   } 
	 & & \\ 
 

    \theutterance \stepcounter{utterance}  

    & & & \multicolumn{2}{p{0.3\linewidth}}{\cellcolor[rgb]{0.95,0.95,0.95}{%
	\makecell[{{p{\linewidth}}}]{% 
	  \tt {\tiny [GM$|$GM]}  
	 * Main Score: 0.00\\ \tt \\ \tt -------\\ \tt * Distance Score: 0.78\\ \tt * Consistency Score: 0.67\\ \tt * Coverage Score: 0.75\\ \tt * Penalty Score: 0.00\\ \tt * Alternative Penalty Score: 0.00\\ \tt * Alternative Main Score: 0.00\\ \tt \\ \tt -------\\ \tt * Shifts: 3.00\\ \tt * Max Shifts: 4.00\\ \tt * Min Shifts: 2.00\\ \tt * End Distance Sum: 2.24\\ \tt * Init Distance Sum: 8.73\\ \tt * Expected Distance Sum: 12.57\\ \tt * Penalties: 9.00\\ \tt * Max Penalties: 8.00\\ \tt * Rounds: 7.00\\ \tt * Max Rounds: 12.00\\ \tt * Object Count: 3.00\\ \tt  
	  } 
	   } 
	   } 
	 & & \\ 
 

\end{supertabular}
}

\end{document}
